\documentclass[10pt,a4paper]{amsart}
\usepackage{amsmath,amssymb,amsthm,graphicx}
\usepackage{amsmath}
\usepackage{amssymb}
\usepackage{hyperref}
% Minimal class-compatibility preamble for a reader excerpt of
%   Juan Esteban Rodriguez Camargo, "An integral model of the perfectoid modular curve",
%   Journal de l'Ecole polytechnique -- Mathematiques 8 (2021), DOI 10.5802/jep.170.
%
% The block below is the author's own macro and theorem declaration block of the pinned
% source file, copied mechanically and unchanged (source lines 6-91), followed by the
% abbreviations that the publisher class cedram supplies and that the retained author text
% uses (same meanings as in that class). No mathematical macro, statement, symbol or
% assumption is introduced, changed or removed by this file.
\usepackage{xspace}

\usepackage[scr=rsfs,cal=euler]{mathalfa}
\multlinegap0pt

\newenvironment{enumeratei}
{\bgroup\def\theenumi{\roman{enumi}}\def\theenumii{\arabic{enumii}}\begin{enumerate}}
{\end{enumerate}\egroup}

\newcommand\bw[1]{\mathbin{\hbox{$\mathsf{\Lambda}$}^{\!#1}}}
\newcommand{\quand}{\quad\text{and}\quad}
\newcommand{\sbullet}{{\scriptscriptstyle\bullet}}

\newcommand\mto{\mathchoice{\longmapsto}{\mapsto}{\mapsto}{\mapsto}}
\newcommand\hto{\mathrel{\lhook\joinrel\to}}
\newcommand\To[1]{\mathchoice{\xrightarrow{\textstyle\kern4pt#1\kern3pt}}{\stackrel{#1}{\to}}{}{}}
\let\ra\rightarrow

\let\ts\textstyle
\newcommand{\lcr}{[\![}
\newcommand{\rcr}{]\!]}
\newcommand{\lpr}{(\!(}
\newcommand{\rpr}{)\!)}

\newcommand{\rae}{\mathrm{ae}}
\newcommand{\an}{\mathrm{an}}
\newcommand{\rb}{\mathrm{b}}
\newcommand{\rc}{\mathrm{c}}
\newcommand{\qc}{\mathrm{qc}}
\newcommand{\cusp}{\mathrm{cusp}}
\newcommand{\cusps}{\mathrm{cusps}}
\newcommand{\cyc}{\mathrm{cyc}}
\newcommand{\et}{\mathrm{\acute{e}t}}
\newcommand{\fTd}{\mathrm{fTd}}
\newcommand{\pr}{\mathrm{pr}}
\newcommand{\proet}{\mathrm{pro}\textup{-}\et}
\newcommand{\sm}{\mathrm{sm}}
\newcommand{\univ}{\mathrm{univ}}

\newcommand{\wotimes}{\mathbin{\widehat\otimes}}

\usepackage[all,cmtip]{xy}\let\labelstyle\textstyle
\usepackage{tikz-cd}

\numberwithin{equation}{section}

\newtheorem{prop}{Proposition}[section]
\newtheorem{lem}[prop]{Lemma}
\newtheorem{theo}[prop]{Theorem}
\newtheorem{cor}[prop]{Corollary}

\theoremstyle{definition}
\newtheorem*{claim}{Claim}
\newtheorem{exam}[prop]{Example}
\newtheorem{definition}[prop]{Definition}
\newtheorem{remark}[prop]{Remark}

\newcommand{\bb}[1]{\mathbb{#1}}
\newcommand{\n}[1]{\mathcal{#1}}
\newcommand{\f}[1]{\mathfrak{#1}}
\newcommand{\s}[1]{\mathscr{#1}}

\DeclareMathOperator{\GL}{GL}
\DeclareMathOperator{\Tr}{Tr}
\DeclareMathOperator{\id}{id}
\DeclareMathOperator{\SL}{SL}
\DeclareMathOperator{\R}{R}
\DeclareMathOperator{\D}{D}
\DeclareMathOperator{\N}{N}

\DeclareMathOperator{\Hom}{Hom}
\DeclareMathOperator{\dotimes}{\otimes^{\mathrm{L}}}
\DeclareMathOperator{\coker}{coker}

\DeclareMathOperator{\Gal}{Gal}
\DeclareMathOperator{\Spf}{Spf}
\DeclareMathOperator{\Spec}{Spec}
\DeclareMathOperator{\Ell}{Ell}
\DeclareMathOperator{\Spa}{Spa}
\DeclareMathOperator{\Cusps}{Cusps}
\DeclareMathOperator{\Tate}{Tate}
\DeclareMathOperator{\sHom}{\s{H}\!\mathit{om}}
\DeclareMathOperator{\HomSurj}{HomSurj}
\DeclareMathOperator{\Art}{Art}
\DeclareMathOperator{\Intperf}{\f{I}nt\textup{-}\f{P}erf}
\DeclareMathOperator{\Perf}{\s{P}erf}
\newcommand{\mc}[1]{X(Np^{#1})}
\newcommand{\fmc}[1]{\mathfrak{X}(Np^{#1})}

% Publisher-class abbreviations (cedram) used by the retained author text.
\providecommand{\loccit}{\emph{loc. cit.}}
\providecommand{\bibliofont}{\small}
\providecommand{\backmatter}{}
\providecommand{\skpt}{\smallskip}
\providecommand{\smaller}{\small}
\newcommand{\ie}{i.e.,\xspace}
\newcommand{\eg}{e.g.,\xspace}
\newcommand{\etc}{etc.\xspace}
\newcommand{\cf}{cf.\xspace}
\newcommand{\resp}{resp.\xspace}
\title{The pi-adic completion of a localization of an integral perfectoid ring is again integral perfectoid}
\author{Juan Esteban Rodríguez Camargo}
\date{doi:10.5802/jep.170; CC BY 4.0}
\begin{document}
\maketitle
\noindent\textit{Editorial note.} All delivered mathematical text is the author's own, copied from the pinned publisher source sources/algebra-fl03-an-integral-model-of-the-perfectoi/pinned-original-source.tex; no statement or proof step was rewritten, reordered or removed.\par\medskip
\noindent\textit{Editorial note.} Preamble provenance: the retained preamble is the source's own retained-preamble.tex, which carries the author's macro and theorem declaration block verbatim (its own header records which lines of the pinned file it copies); the wrapper class supplies the baseline packages that block assumes. Nothing was added, no macro was invented, and the excerpt serves only as preamble, never as a span coordinate.\par\medskip
\noindent\textit{Editorial note.} Metadata provenance: source\_id and source\_version were taken from sources/algebra-fl03-an-integral-model-of-the-perfectoi/license-observation.json; the title from SOURCE.json; the author list from the source metadata record; the source URL from the first non-empty of source\_url, source\_urls[0] and the evidence's fixed\_version\_url/license\_url. The licence evidence file used by this card is sources/algebra-fl03-an-integral-model-of-the-perfectoi/license-observation.json, recorded in license.evidence\_path so that the binding target is explicit.\par\medskip
\noindent\textit{Editorial note.} External references are navigation pointers only; every deliberately kept step of the author's argument stays inside the counted spans. The displayed numbers of those pointers are the author's own numbering under the author's \textbackslash{}newtheorem declarations (section.index as declared), computed mechanically by the lane's scanner.\par\medskip
\noindent\textit{Editorial note.} The counted proof contains no reference to any labelled item and no bibliography citation, so the card needs neither a navigation stub nor a bibliography excerpt.\par\medskip
\noindent\textit{Editorial note.} This is the first card this lane authors from this source, taken from the newly approved batch of sources; it was selected as a zero-stub, zero-citation candidate to keep the publishable queue fed.\par\medskip
\noindent\textit{Allowed standard prerequisites (editorial).}\begin{itemize}
\item The definition of an integral perfectoid ring with pseudo-uniformizer pi and the attached formal scheme. Set up in the paragraph preceding the counted statement (source line 467). Fixes the two properties that have to be verified for the completed localization: pi-adic completeness with pi a non-zero-divisor, and the Frobenius being surjective on R/pi.
\item The short exact sequences R/pi\textasciicircum{}n -> R/pi\textasciicircum{}\{n+k\} -> R/pi\textasciicircum{}k -> 0 coming from the fact that pi is not a zero divisor. Used inside the counted proof at lines 475-478. Supplies the sequences that are localized and then passed to the inverse limit.
\end{itemize}\medskip
\begin{lem}
\label{lemLocalizationperf}
Let $f\in R$. Then $R\langle f^{-1} \rangle= \varprojlim_n R/\pi^n [f^{-1}]$ is an integral perfectoid ring.
\end{lem}

\begin{proof}
Let $n,k\geq 0$, as $\pi$ is not a zero divisor we have a short exact sequence
\[
0\to R/\pi^n \To{\pi^k} R/\pi^{n+k}\to R/\pi^k \to 0.
\]
Localizing at $f$ and taking inverse limits on $n$ we obtain
\[
0\to R\langle f^{-1}\rangle\To{\pi^k} R\langle f^{-1} \rangle\to R/\pi^{k} [f^{-1}]\to 0.
\]
Then $R\langle f^{-1} \rangle$ is $\pi$-adically complete and $\pi$ is not a zero-divisor. On the other hand, localizing at $f$ the Frobenius map
$\varphi: R/\pi \To{\sim} R/\pi^p$ one gets
\[
\varphi: R/\pi[f^{-1}]\To{\sim} R/\pi^{p}[f^{-p}]= R/\pi^p [f^{-1}],
\]
which proves that $R\langle f^{-1}\rangle$ is an integral perfectoid ring.
\end{proof}

\begin{thebibliography}{99}
\bibitem{hartshorne1966residues} R. Hartshorne - Residues and duality , Lect. Notes in Math. , vol. 20 , Springer-Verlag, Berlin-New York, 1966
\bibitem{mumford1974abelian} D. Mumford - Abelian varieties , TIFR Studies in Mathematics , vol. 5 , Tata Institute of Fundamental Research, Bombay, 2008
\end{thebibliography}
\end{document}
